原始手稿提出了一个有用的核心判断：人们之所以经常无法启动那些明知有益的行动，并不只是因为他们缺少价值判断，而是因为他们被嵌在某个既有的行为框架里，这个框架本身的回报结构、认知惯性与所处环境都会共同抗拒转变。这个判断仍然是整个项目的中心。这一稿改变的，是它的教学雄心。

这一版写成的是教材原型，而不再是一篇较短的概念性文章。教材需要的，不只是正确的直觉。它还需要有序的定义、稳定的记号、逐章推进的结构、展开后的案例，以及一套能让读者检验这个框架究竟是否真的解释了行为，而不只是听起来有启发性的练习层。